\documentclass{beamer}
\usepackage[utf8]{inputenc}
\usepackage[spanish]{babel}
\usetheme{Boadilla}
\usecolortheme{beaver}

\title{Energías renovables en Colombia}
\subtitle{Proyecto de aula --- Ciencias Ambientales}
\author{Valentina Ríos \and Mateo Cárdenas}
\institute{Colegio San Bartolomé}
\date{\today}

\begin{document}

\begin{frame}
  \titlepage
\end{frame}

\begin{frame}{Contenido}
  \tableofcontents
\end{frame}

\section{Situación actual}
\begin{frame}{Matriz energética}
  \begin{itemize}
    \item Cerca del 65\,\% de la electricidad proviene de hidroeléctricas.
    \item La dependencia del agua genera riesgos en épocas de sequía.
    \item La energía solar y eólica crecen rápidamente en La Guajira.
  \end{itemize}
\end{frame}

\section{Alternativas}
\begin{frame}{Fuentes renovables no convencionales}
  \begin{columns}[T]
    \begin{column}{0.48\textwidth}
      \textbf{Solar}
      \begin{itemize}
        \item Alta radiación todo el año
        \item Costos en descenso
      \end{itemize}
    \end{column}
    \begin{column}{0.48\textwidth}
      \textbf{Eólica}
      \begin{itemize}
        \item Vientos constantes en el Caribe
        \item Retos de consulta previa
      \end{itemize}
    \end{column}
  \end{columns}
\end{frame}

\section{Conclusiones}
\begin{frame}{Conclusiones}
  \begin{block}{Idea principal}
    Diversificar la matriz energética reduce la vulnerabilidad frente al cambio
    climático.
  \end{block}
  \pause
  ¡Gracias por su atención!
\end{frame}

\end{document}
